\documentclass{beamer}

\usepackage{amsmath}
\usepackage{amsfonts}
\usepackage{amssymb}
\usepackage{textpos}
\usepackage[orientation=landscape,size=custom,width=1900,height=1100,scale=1.2]{beamerposter}
\usepackage{ragged2e}
\usepackage{graphicx}

\usepackage{xcolor}
\definecolor{darkorchid}{HTML}{9932CC}
\definecolor{berlinred}{rgb}{0.82, 0.1, 0.26}

\usepackage{array}
\usepackage{booktabs}
\usepackage{bbding}

\usepackage{caption}
\usepackage{subcaption}

% Global default (e.g., tables) stays centered
\captionsetup{justification=centering}
% Override only for figures
\captionsetup[figure]{justification=justified,singlelinecheck=false}
\captionsetup[subfigure]{justification=justified,singlelinecheck=false}

\usetheme{Berlin}
\setbeamertemplate{headline}{}
\setbeamertemplate{footline}{}

% Redefine block environment background to white
\setbeamercolor{block body}{bg=white}
\setbeamercolor{block title}{bg=purple}
\setbeamercolor{title}{bg=purple}

\begin{document}

\title[]{\Huge Container-based High-Performance Computing Approaches for Simulating Large-Scale Biological Neural Networks  \\[-0.1 in]}
\author[]{\LARGE \textbf{Oliver James} and \textbf{Sungho Hong$\dagger$ } \\[-0.1 in] }
\institute[]{\Large Center for Cognition and Sociality, Institute for Basic Science;~$\dagger$ sunghohong@ibs.re.kr }
\date{}

\begin{frame}
\vspace{-0.4 in}
\maketitle

\begin{columns}[t,totalwidth=\paperwidth]

    % =================== COLUMN 1 ===================
    \begin{column}{0.32\paperwidth}

        \begin{block}{\LARGE Introduction}
        \justifying
        {\large
        1. Modeling \textbf{large-scale, biophysically detailed neural networks} increasingly relies on \textbf{high-performance computing (HPC)}. \\[0.5 cm]

        2. Deploying such simulations across diverse HPC environments is challenging due to \textbf{software incompatibilities} and \textbf{complex dependencies}. \\[0.5 cm]

        3. We address this using \textbf{Singularity containers} that encapsulate the entire environment, enabling portability and reproducibility without root privileges.\\[0.5 cm]

        4. Practical demonstrations: benchmarking HPC environments with \textbf{cerebellar mossy fiber$-$granule cell network} simulations.
        }
        \end{block}

        \begin{block}{\LARGE Containerization for Scientific Computing}
        \justifying
        {\large
        A container is a lightweight, standalone, virtualized environment that includes all dependencies required to run large-scale scientific models. \\[1 cm]
        }
        \centering
        \includegraphics[scale=1.8]{output1.eps}
        \end{block}

    \end{column}

    % =================== COLUMN 2 ===================
    \begin{column}{0.32\paperwidth}

        \begin{block}{\LARGE Our Computational Framework}
        \centering
        \includegraphics[scale=1.6]{output.eps}
        \end{block}

        \begin{block}{\LARGE Benchmarking Data and Network: Cerebellar Mossy Fiber $\rightarrow$ Granule Cell}
        \centering
        \includegraphics[scale=1.7]{CC.eps}
        \end{block}

        \justifying
        {\large We run multiple copies ($n=280$) of the network model simultaneously to simulate extreme workloads. \\[3 mm]

        This allows rigorous assessment of GPU scalability and performance limits.\\[5 mm]}

        \begin{table}[htbp]
        \centering
        \begin{minipage}{0.5\linewidth}
        \centering
        \caption{Network Configuration}
        \begin{tabular}{c@{\hspace{2cm}}c}
        \toprule
        \textbf{Network elements} & \textbf{Count (M)} \\
        \midrule
        \textcolor{green}{MF} & 0.299 \\
        \textcolor{gray}{GRC} & 1.099 \\
        Compartments & 3.297 \\
        Synapses & 3.468 \\
        \bottomrule
        \end{tabular}
        \end{minipage}%
        \begin{minipage}{0.45\linewidth}
        \centering
        \caption{Simulation Setup}
        \begin{tabular}{l@{\hspace{1.5cm}}c}
        \toprule
        Simulation time & 1\,s \\
        Software & CoreNEURON \\
        \bottomrule
        \end{tabular}
        \end{minipage}
        \end{table}

    \end{column}

    % =================== COLUMN 3 ===================
    \begin{column}{0.32\paperwidth}

        \begin{block}{\LARGE Benchmarking HPC Environments}
        \begin{table}[htbp]
        \centering
        \caption{Performance in CPU parallel computing (30-task MPI)}
        \setlength{\tabcolsep}{25pt}
        \begin{tabular}{c c c  c }
        \hline
        CPU & Cores & Solver Time (s) & Simulator \\
        \hline
        Intel Xeon & 30 & 1670  & CoreNEURON-CPU \\
        \hline
        \end{tabular}
        \end{table}

        \begin{figure}[htbp]
            \centering
            \begin{subfigure}[t]{0.48\textwidth}
                \centering
                \includegraphics[width=\textwidth]{Turing.eps}
                \caption{TITAN RTX: Solver time vs. GPU count}
                \label{fig:turing}
            \end{subfigure}
            \hfill
            \begin{subfigure}[t]{0.48\textwidth}
                \centering
                \includegraphics[width=\textwidth]{HPCcompare.eps}
                \caption{Comparison across HPC environments}
                \label{fig:hpccompare}
            \end{subfigure}
        \end{figure}

        \justifying
        {\large TITAN RTXs outperform 30-core CPUs by up to 9.8×, with performance improving as GPU count increases.}
        \end{block}

        \begin{block}{\LARGE Future Works}
        \justifying
        \HandRight~ \textbf{Develop a biologically realistic cerebellar circuit model} (granule, Golgi, Purkinje, interneurons). \\[5 mm]
        \HandRight~ \textbf{Benchmark across institutions} to assess scalability and performance of HPC systems.
        \end{block}

        {
        \setbeamercolor{myblock}{bg=berlinred!100, fg=white}
        \begin{beamercolorbox}[wd=\columnwidth, sep=1.5ex]{myblock}
        {\Large Acknowledgment}
        \end{beamercolorbox}}
        \justifying
        {\large This work was supported by IBS-R001-D2-2025-a004D5}

        \begin{flushright}
        
\includegraphics[scale=1.0]{H.eps}
        \end{flushright}

    \end{column}

\end{columns}

\end{frame}
\end{document}
